\documentclass[11pt]{article}
\usepackage{amsmath,amssymb}
\title{Small rejection certificates for several timed duties}
\date{25 November 2026}
\begin{document}
\maketitle

\paragraph{Interface.}
Extend the fixed-mode interval decision to a vector of \(d\) event
times, with common anchor \(t_0=0\). The receiver chooses the mode
and a vector before it can distinguish the admitted histories.
For each fixed mode \(m\) and history \(h\), assume the complete
admissible schedule is given by a finite conjunction of integer
difference bounds \(t_j-t_i\leq b\), together with the common
ambient bounds. The choice of mode must be the same for all
histories at that cut. The event correspondence across sources
must be part of the interface: timestamps referring to issue,
actuation, and delivery cannot be interchanged. This treats fixed
internal execution order and a finite decision cut; arbitrary
disjunctions or observation-dependent future execution are outside
its completeness claim.

\paragraph{Exact decision and witness size.}
Combine the bounds in each mode, retaining the smallest bound
for each oriented pair and one actual source history attaining it.
This preserves the feasible schedules. Regard each bound as a
weighted arc \(i\to j\). Feasibility is equivalent to absence of a
negative directed cycle. Necessity is telescoping along the cycle.
For sufficiency, a zero-weight super-source reaches every event.
If every cycle has nonnegative weight, shortest path weights are
attained on simple paths and satisfy every arc inequality.
They are integers for integer bounds, and subtracting the anchor
value makes \(t_0=0\) without changing differences.
This is a standard simple temporal network calculation.

In every rejected mode a negative simple cycle uses at most \(d+1\)
arcs. Consequently across \(r\) modes, at most \(r(d+1)\)
histories suffice to refute a shared schedule. A history can be
reused by several arcs; common bounds consume no history witnesses.
The count is sharp. For each mode assign \(d+1\) distinct histories
one arc each from the cycle
\[
  0\to1\to\cdots\to d\to0
\]
with all bounds zero except the last bound \(-1\).
All other modes receive only ambient bounds from these histories,
and all event times are in \([0,1]\).
Removing the closing arc permits every event at zero; removing the
arc \(i\to i+1\) permits event times zero through \(i\), then one.
Omitting any one history restores feasibility in its mode. For
one event variable, the two-arc case is precisely the interval
bound, though applications may admit stronger common bounds.

\paragraph{Physical read cut and revalidation.}
The two-report mechanical-flag example gives an explicit source
projection. Let \(i,\ell,h,b\) denote actual issue, listen, handoff,
and final delivery times. Prior work requires \(i\geq u\);
the first report forces \(\ell\geq2\); actual flag readiness requires
\(\ell-i\geq L\); listening and movement require \(h-\ell\geq2\);
the custody transfer and return require \(b-h\geq3\).
With final deadline seven and frozen read at two, \(u=0,L=2\)
admits \((0,2,4,7)\). With \(u=1,L=2\), the issue and latency
constraints imply \(\ell\geq3\), directly refuting the frozen read.
Even allowing re-timing without changing the old deadline gives a
different negative cycle. Deadline eight with a newly timed receiver
admits \((1,3,5,8)\). The timings reflect the stated mode of the
physical reference; they are not a new whole-history planner.

After a source replacement, a retained negative cycle remains a
valid witness if every arc has a certified bound on the same pair
of events that is no weaker, and those facts are admitted at the
decision cut. All other histories may change. If a witnessing arc
is forgotten, reversed in meaning, or weakened, the old proof
no longer gives that verdict. If a previously impossible mode set
is enlarged, a witness for each old mode does not refute the new one.
This is a proof reuse condition, not a claim that discovering or
validating source edits is free.

\paragraph{Verification and comparison.}
Lean checks telescoping, the negative-cycle obstruction and the
arc-by-arc replacement rule in \texttt{TimedEventCertificates.lean}.
Bellman--Ford feasibility and the sharp history count are proved
above, not verified as executable Lean code. The replay compares
with exhaustive integer timeline choices on 2,400 model instances,
tests the tight family, and reproduces both physical read-cut cases.
Direct simple temporal network solving matches the guarantee.
Per-edge indexing may reduce repeated reads only when the source
provides a complete witnessed update stream. We do not claim a
new graph theorem, speedup over a factored matched implementation,
or acceptance of R2 or the integrated stack.
\end{document}
